\documentclass[10pt]{article}
\usepackage[T1]{fontenc}
\usepackage[utf8]{inputenc}
\usepackage[french]{babel}
\usepackage{amsmath,amssymb,mathrsfs,mathpazo}
\usepackage{geometry}
\geometry{a4paper,margin=2cm}
\usepackage{graphicx,array,tabularx,tabularray,colortbl,multirow,multicol,booktabs,enumitem,pifont,ulem,fancybox,diagbox,cancel}
\usepackage{pstricks,pstricks-add,pst-plot,pst-tree,pst-node,pst-text,pst-func,pst-math,pst-bspline,pst-3dplot,pst-eucl,pst-solides3d}
\let\pscalculate\undefined
\usepackage{pst-bezier}
\usepackage{tikz,pgfplots,tkz-tab}
\usetikzlibrary{arrows,arrows.meta,patterns,calc,positioning,trees,decorations.markings,shapes.geometric,shapes.misc,shapes.arrows,intersections,decorations.pathmorphing,decorations.pathreplacing,backgrounds,shadows,babel}
\pgfplotsset{compat=1.18}
\usepgfplotslibrary{fillbetween}
\usepackage[np]{numprint}
\usepackage{listings,makecell,pas-tableur,twoopt}
\let\barre\undefined
\usepackage[pstricks]{tablvar}
\definecolor{darkgreen}{RGB}{0,100,0}
\providecommand{\red}{\color{red}}
\providecommand{\blue}{\color{blue}}
\providecommand{\green}{\color{green}}
\providecommand{\black}{\color{black}}
\providecommand{\white}{\color{white}}
\providecommand{\gray}{\color{gray}}
\providecommand{\cyan}{\color{cyan}}
\providecommand{\magenta}{\color{magenta}}
\providecommand{\yellow}{\color{yellow}}
\providecommand{\violet}{\color{violet}}
\providecommand{\orange}{\color{orange}}
\providecommand{\euro}{\texteuro}
\providecommand{\EUR}{\texteuro}
\providecommand{\ieme}{\textsuperscript{e}}
\providecommand{\degres}{\ensuremath{^\circ}}
\providecommand{\degree}{\ensuremath{^\circ}}
\providecommand{\vv}[1]{\overrightarrow{#1}}
\providecommand{\Coord}[1]{#1}
\providecommand{\decosix}{}
\providecommand{\starredbullet}{\ensuremath{\bullet}}
\pagestyle{empty}
\setlength{\parindent}{0pt}
\usepackage{adjustbox}\DeclareUnicodeCharacter{2236}{:}\DeclareUnicodeCharacter{2190}{\ensuremath{\leftarrow}}
\providecommand{\R}{\mathbb{R}}
\providecommand{\N}{\mathbb{N}}
\providecommand{\Z}{\mathbb{Z}}
\providecommand{\Q}{\mathbb{Q}}
\providecommand{\D}{\mathbb{D}}
\providecommand{\C}{\mathbb{C}}
\providecommand{\e}{{\mathrm{e}}}
\providecommand{\vect}[1]{\overrightarrow{#1}}
\providecommand{\vectt}[1]{\overrightarrow{\mathrm{#1}}}
\providecommand{\vvt}[1]{\overrightarrow{\mathrm{#1}}}
\providecommand{\np}[2][]{{#2~#1}}
\providecommand{\vv}[1]{{\overrightarrow{#1}}}
\renewcommand{\d}{{\mathrm{d}}}
\providecommand{\degres}{{^\circ}}
\providecommand{\degre}{{^\circ}}
\providecommand{\blue}{\color{blue}}
\providecommand{\red}{\color{red}}
\providecommand{\enskip}{\quad}
\providecommand{\llbracket}{\lbrack\!\lbrack}
\providecommand{\rrbracket}{\rbrack\!\rbrack}
\providecommand{\euro}{{\text{€}}}
\providecommand{\no}{{\text{n°}}}
\providecommand{\interval}[2]{{[#1;#2]}}
\providecommand{\intervFO}[2]{{[#1;#2[}}
\providecommand{\dotfill}{\ldots}

\providecommand{\ds}{\displaystyle}
\providecommand{\textding}[1]{#1}
\providecommand{\Oij}{\(\mathrm{O};\vec{\imath},\vec{\jmath}\)}
\providecommand{\Oijk}{\(\mathrm{O};\vec{\imath},\vec{\jmath},\vec{k}\)}
\providecommand{\Ouv}{\(\mathrm{O};\vec{u},\vec{v}\)}
\providecommand{\intervFF}[2]{\ensuremath{[#1;#2]}}
\renewcommand{\euro}{\ifmmode\text{\texteuro}\else\texteuro\fi}
\DeclareUnicodeCharacter{00B0}{\ifmmode^{\circ}\else\textdegree\fi}
\DeclareUnicodeCharacter{20AC}{\ifmmode\text{\texteuro}\else\texteuro\fi}
% Énoncés : Éducation nationale. Transcription et corrigés : APMEP, Denis Vergès et contributeurs. Découpage et classement : CoopMaths, Sébastien Lozano, Éric Elter et contributeurs. Adaptation de la présentation : Maths Explorer.
% CC BY-SA : https://coopmaths.fr/www/about/
% https://coopmaths.fr/alea/static/eam/2026/tex/eam_technologique_2026_06_etrangers_2_cor.tex
\begin{document}
\noindent\textit{APMEP / CoopMaths. CC BY-SA. Présentation adaptée par Maths Explorer.}\par\medskip
\begin{enumerate}
\item Par lecture graphique :
\begin{enumerate}
\item Le point de la courbe d'abscisse 3 a pour ordonnée $-2$. Donc : $f(3)=-2$.

\item Le nombre dérivé $f'(3)$ est le coefficient directeur de la tangente $T$ au point d'abscisse 3.

La droite $T$ passe par le point $(3\,;\,-2)$ et, par lecture graphique, par le point $(4\,;\,2)$. 

Son coefficient directeur est donc : $\dfrac{2-(-2)}{4-3}=\dfrac{4}{1}=4$.

Ainsi : $f'(3)=4$.

\end{enumerate}

\item On considère désormais la fonction définie sur $[0\,;\,4]$ par : $f(x)=-x^3+9x^2-24x+16$.

\begin{enumerate}
\item On calcule : $f(0)=-0^3+9\times 0^2-24\times 0+16=16$.

\item Pour tout réel $x$ de l'intervalle $[0\,;\,4]$ :
\[f'(x)=-1\times3x^2+9\times x-24+0=-3x^2+18x-24.\]

\item On développe :
\[-3(x-2)(x-4)=-3(x^2-4x-2x+8)=-3(x^2-6x+8)=-3x^2+18x-24.\]
On retrouve bien l'expression de $f'(x)$ donnée.
\end{enumerate}
\end{enumerate}

\begin{enumerate}
\item[]
\begin{enumerate}[start=4]
\item On utilise la forme factorisée : $f'(x)=-3(x-2)(x-4)$.

Le facteur $x-2$ est nul en 2, le facteur $x-4$ est nul en 4.

Calculons les valeurs utiles :
\[\text{On a déjà : } f(0)=16,\]
\[f(2)=-2^3+9\times 2^2-24\times 2+16=-8+36-48+16=-4,\]
\[f(4)=-4^3+9\times 4^2-24\times 4+16=-64+144-96+16=0.\]

On obtient :

\begin{center}
\begin{tikzpicture}
\tkzTabInit[lgt=2.8,espcl=2.8]{$x$/1,$-3$/1,$x-2$/1,$x-4$/1,$f'(x)$/1,$f(x)$/2}{$0$,$2$,$4$}
\tkzTabLine{,-,t,-,t}
\tkzTabLine{,-,z,+,t}
\tkzTabLine{,-,t,-,z}
\tkzTabLine{,-,z,+,z}
\tkzTabVar{+/$16$,-/$-4$,+/$0$}
\end{tikzpicture}
\end{center}

\item D'après le tableau de variation, le minimum de la fonction $f$ sur l'intervalle $[0\,;\,4]$ est $y=-4$, il est atteint pour $x=2$.

\end{enumerate}
\end{enumerate}

\end{document}
